Our algorithm for RIT uses non-determinism to guess a "good" prime.  It is natural to wonder whether such primes can instead be randomly sampled.
Recall we require that the prime $p$ have polynomial bit-length in the size of the input, and that  the congruences $x^{d_i} \equiv a_i \pmod{p}$ are solvable in $\FF_p$. By Chebotarev's density theorem, roughly speaking, the density of such good primes is
$\frac{1}{|\Gal(K/\QQ)|}$.
Since the size of the Galois group of $K$ over $\QQ$ is exponential in the size of the input, good primes do not have sufficient density in order to directly be chosen randomly. 
The density remains insufficient even if the exponents~$d_i$ are prime numbers written in unary.
 

In what follows, we show that   $2$-RIT is in {\coRP} under GRH, and  in {\coNP} unconditionally. Recall that  $2$-RIT  is the identity testing problem for an algebraic circuit~$C$ evaluated on square-roots
$\sqrt{a_1},\ldots,\sqrt{a_k}$ for 
$k$ rational odd primes $a_1,\ldots,a_k$.
The proofs from \Cref{sec:correctness} ensure that the finite field computation in our algorithm is sound; it remains to show how to choose a completely split prime~$p$ and determine the solutions to the equations $x^{2} \equiv a_i \pmod{p}$  in~$\FF_p$.


As noted above, 
the natural density of primes is not sufficient even for  $2$-RIT, 
however, we show that there is an arithmetic progression with a good density of primes, and that all primes in this progression are good.

Below, we state known effective bounds on the density of primes in an arithmetic progression, in particular, the following estimates, which have been shown in \cite[Chapter 20, page 125]{davenport} and \cite[Corollary 18.8]{IKbook}, respectively:

\begin{theorem}\label{th:prime_density}
	Given $a\in\ZZ_n^*$, write $\pi_{n,a}(x)$ for the number of primes less than $x$ that are congruent to $a$ modulo $n$. Then under GRH, there is an absolute constant $c>0$ such that
	\[ \pi_{n,a}(x) \geq \frac{x}{\varphi(n)\log x} - c x^{1/2}\log x.\]
	Unconditionally, there exist effective positive constants $c_1$ and $c_2$, such that for all $n < c_1 x^{c_1}$,
	\[ \pi_{n,a}(x) \geq \frac{c_2 x}{\varphi(n)x^{1/2}\log x}. \]
\end{theorem}

What remains to be understood is how to construct an arithmetic progression such that for every prime $p$ appearing in it, $\FF_p$ contains a representation $\overline{\alpha}_1,\ldots,\overline{\alpha}_k\in\FF_p$ of the square-root input $\sqrt{a_1},\ldots,\sqrt{a_k}$. As it turns out, this only
requires some classical results on quadratic reciprocity,
which we recall now.

Let $p$ be an odd prime number. An integer $a$ is said to be a \emph{quadratic residue} modulo $p$ if it is congruent to a perfect square modulo $p$, i.e., if there exists an integer $x$ such that $x^2 \equiv a \mod p$. The \emph{Legendre symbol} is a function of $a$ and $p$ taking values in~$\{1,-1,0\}$, that is
\[\left( \frac{a}{p} \right) =\begin{cases}
1 & \text{if $a$ is a quadratic residue mod $p$ and $a \not\equiv 0 \pmod{p}$},\\
-1 &\text{if $a$ is a non-quadratic residue mod $p$},\\
0 & \text{if $a\equiv 0 \pmod{p}$.} 	
\end{cases}
\]
  Its explicit definition is as follows:
\[ \left( \frac{a}{p} \right) =a^{\frac{p-1}{2}} \pmod{p}. \]
Furthermore, given  odd primes $p$ and $q$, the \emph{Law of quadratic reciprocity} states:
\[ \left( \frac{p}{q} \right)\left( \frac{q}{p} \right) = (-1)^{\frac{p-1}{2}\frac{q-1}{2}}.\]
We observe that in the case where $p\in 4\NN+1$, then
\[ \left( \frac{p}{q} \right)\left( \frac{q}{p} \right) = 1 \; \; \Leftrightarrow \; \; \left( \frac{p}{q} \right) = \left( \frac{q}{p} \right) = \pm 1.\]
In other words, $p$ is a quadratic residue modulo $q$ if and only if $q$ is a quadratic residue modulo~$p$, when either $p$ or $q \equiv 1 \pmod{4}$.



\begin{figure*}[t]
    \centering
    \begin{tabularx}{0.82\textwidth}{ p{0.1\textwidth} p{0.72\textwidth} }
   \rowcolor{Gray}
    \multicolumn{2}{c}{\textbf{Radical Identity Testing for square root inputs}} \\
    \textbf{Input:} & Algebraic circuit $C$ of size  at most $s$
    and a list of $k$ odd primes $a_1, \ldots, a_k$ 
    of magnitude at most $2^s$.
    \\
    \rowcolor{Gray}
    \textbf{Output:} & Whether $f(\sqrt{a_1},\ldots,\sqrt{a_k})=0$ for the polynomial $f(x_1,\ldots,x_k)$ computed by $C$.
    \\ 
    \textbf{Step 1:} &
            Compute $b$ such that $b +1 \equiv 5 \pmod{8}$, and $b +1 \equiv 1\pmod{a_i}$ for all $i$.
    \\  
    \textbf{Step 2:} & 
    		Pick $p$ uniformly at random from the set $S(a_1,\ldots,a_k)$ defined in~(\ref{eq:SA}).
     \\ 
    \textbf{Step 3:} &
            Compute $\overline{\alpha}_1,\ldots,\overline{\alpha}_k \in\FF_p$ such that $\overline{\alpha}_i^{2}  \equiv a_i \pmod{p}$ as described in
            \Cref{eq:pocklington}.
     \\  
    \textbf{Step 4:} &
            Output 'Zero' if $\overline{f}(\overline{\alpha}_1,\ldots,\overline{\alpha}_k) = 0$, where $\bar{f}$ is the reduction of $f$ modulo $p$; and 'Non-zero' otherwise.
    \\ \hline
    \end{tabularx}
    \caption{Our randomised polynomial time
   algorithm for the complement of  $2$-RIT.
    }
    \label{fig:corp_squareRIT}
\end{figure*}


In what follows, we apply the law of quadratic reciprocity in order to choose the right field $\FF_p$ for deciding $2$-RIT. Recall that intuitively, we are looking for a prime $p$ such that $x^2-a_i$ has a solution in $\FF_p$ for all $i$, that is, the $a_i$ is a quadratic residue modulo $p$. Since we will be choosing $p$ from an arithmetic progression, we can easily make that progression to be of the shape $4\NN + 1$, that is, to ensure that $p\equiv 1 \pmod{4}$. In that case the $a_i$'s will be quadratic residues modulo $p$ if and only if $p$ is a quadratic residue modulo $a_i$ for all~$i$. In order to ensure that, it suffices to choose $p$ such that $p \equiv 1 \pmod{a_i}$ for all $i$, as $1$ is a perfect square modulo $a_i$ for all $i$, and thus $p$ a quadratic residue modulo $a_i$.

We have just shown that if we choose $p$ such that it satisfies all the above-mentioned congruences, the polynomials $x^{2}-a_i$ all split and have non-zero roots in $\FF_p$. Following Pocklington's algorithm, there is a deterministic way to solve the equations $x^2-a_i$  if $p\equiv 5 \pmod{8}$. 
In particular, writing $p = 8m+5$, the solution of the equation $x^2 \equiv a \pmod{p}$ is given by the following function:
\begin{equation}\label{eq:pocklington}
x =\begin{cases}
\pm a^{m+1} & \text{if $a^{2m+1} \equiv 1 \pmod{p}$,}\\
\pm\frac{y}{2} &\text{if $a^{2m+1} \equiv -1 \pmod{p}$ and} \\
& \text{$y = \pm(4a)^{m+1}$ is even,} \\ 
\pm\frac{p+y}{2} &\text{if $a^{2m+1} \equiv -1 \pmod{p}$ and} \\
& \text{$y = \pm(4a)^{m+1}$ is odd.}
\end{cases}
\end{equation}
See \Cref{appendix:pocklington} for details.
Note that this congruence encompasses the above restriction on $p$ being congruent to 1 modulo~4.


We now show how to construct an arithmetic progression such that all primes $p$ in the progression satisfy the above congruences. Denote by $A=\prod_{i=1,a_i\neq 2}^{k}a_i$ the product of all input radicands $a_i$. Note that $A$ will always be odd as the $a_i$ are odd primes. Let us look at the arithmetic progression
\begin{equation}\label{eq:arithmetic_progression_corp}
	 8A\NN + b + 1,
\end{equation}
where $b$ is a solution of the following system of equations 
\begin{align}\label{eq:bAP}
	b \equiv 4 &\pmod{8} \\
	b \equiv 0 &\pmod{a_i}. \nonumber
\end{align}
Since all the moduli in the equations~(\ref{eq:bAP}) are pairwise coprime, by the Chinese remainder theorem, the  system has a solution.
By the construction above,
 we have an arithmetic progression such that all primes $p$ in the progression are  good primes.
We also ensure that $p$ is such that we can deterministically find the representations $\overline{\alpha}_1,\ldots,\overline{\alpha}_k$ of $\sqrt{a_1},\ldots,\sqrt{a_k}$ in $\FF_p$.
Define by $S(a_1,\ldots,a_k)$ the following set  
\begin{align}\label{eq:SA}\big\{p \leq 2^{5k^3} \mid  &p \in  8A \NN+ b+1  \text{ where } A=\prod_{i=1}^{k}a_i \\
    &\text{ and } b \text{ is a solution of~(\ref{eq:bAP})} \big\}. 	\nonumber
\end{align}
We have the following:

\begin{restatable}{proposition}{probabilityprimep}\label{prop:density_corp}
Let $C$ be an algebraic circuit of size at most $s$, and $a_1, \ldots, a_k$ odd primes of bit-length at most $s$. 
	 Denote by $\alpha$ the algebraic integer obtained by evaluating~$C$ on the~$\sqrt{a_i}$. Suppose that $p$ is chosen uniformly at random from the set~$S(a_1,\ldots,a_k)$ defined in~(\ref{eq:SA}).
	 Then
		\begin{itemize}
			\item[(i)] $p$ is prime with probability at least
			$\frac{1}{6s^3}$ assuming GRH, 
			 and
			\item[(ii)] given that $p$ is prime, the probability that it divides $N(\alpha)$ is at most $2^{-s^3}$	unconditionally.
		\end{itemize}
\end{restatable}



With~\Cref{prop:density_corp} in hand, we can state our algorithm, see \Cref{fig:corp_squareRIT}, and prove its complexity.


\begin{proof}[Proof of \Cref{th:squareRIT-coRP}]
\Cref{fig:corp_squareRIT} presents a 
randomised polynomial time algorithm for the complement of  $2$-RIT  as follows. It is clear that the algorithm runs in polynomial time. Let us now argue its correctness.
	
	First, suppose that $f(\sqrt{a_1},\ldots,\sqrt{a_k}) = 0$. By \Cref{th:RIT-soundess}, we have $\overline{f}(\overline{\alpha}_1,\ldots,\overline{\alpha}_k) = 0$, and hence the output is 'Zero'. Second, suppose that $f(\sqrt{a_1},\ldots,\sqrt{a_k}) \neq 0$. Then the output will be 'Non-Zero' provided that $p$ does not divide $N(f(\sqrt{a_1},\ldots,\sqrt{a_k}))$. By \Cref{prop:density_corp}(ii), the probability that $p$ does not divide $N(f(\sqrt{a_1},\ldots,\sqrt{a_k}))$ is at least $1 - 2^{-s^3}$. 
	Thus, the probability that the algorithm gives the wrong output is at most $2^{-s^3}$.

	 
	It remains to show that  $2$-RIT  is in {\coNP} unconditionally. The idea is to modify the algorithm in \Cref{fig:corp_squareRIT}, replacing randomisation with guessing. \Cref{th:prime_density} shows that \begin{equation}\label{eq:numprimnocon}\pi_{8A,b+1}(2^{3s^3}) > 2^{s^3}\end{equation}
	 for $s$ sufficiently large.
	 It follows that there exists a prime $p\leq 2^{3s^3}$
	 that does not divide $N(f(\sqrt{a_1},\ldots,\sqrt{a_k}))$. The rest of the argument follows as in the proof of Theorem~\ref{th:rit-conp}.
	\end{proof}
	

Note that, in general, in the proof of the theorem above, GRH is required to obtain the {\coRP} bound. 
This is because the unconditional lower bound on density of primes in arithmetic progressions is not strong enough for our purposes: 
The number of primes less than $2^{3s^3}$ which are favourable is just $2^{s^3}$, as computed in \eqref{eq:numprimnocon}. This gives a probability of success at least $2^{-2s^3}$, which is exponentially small in the instance size. In order to get a constant success probability, we have to repeatedly sample and run this algorithm $2^{2s^3}$ times, which yields an exponential time algorithm. However, under GRH, the bound is improved to $\frac{1}{6s^3}$ and polynomially many repetitions suffice for a constant success probability.

	
	
	